% TAB-13: generated by harness/analysis/paper_outputs.py; do not edit (every cell's source is in TAB-13.cells.json)
\begin{tabularx}{\linewidth}{>{\hsize=1.064\hsize\raggedright\hyphenpenalty=10000\exhyphenpenalty=10000\arraybackslash}Xr>{\hsize=1.016\hsize\raggedleft\hyphenpenalty=10000\exhyphenpenalty=10000\arraybackslash}X>{\hsize=1.016\hsize\raggedleft\hyphenpenalty=10000\exhyphenpenalty=10000\arraybackslash}X>{\hsize=1.318\hsize\raggedleft\hyphenpenalty=10000\exhyphenpenalty=10000\arraybackslash}X>{\hsize=0.839\hsize\raggedleft\hyphenpenalty=10000\exhyphenpenalty=10000\arraybackslash}X>{\hsize=0.747\hsize\raggedright\hyphenpenalty=10000\exhyphenpenalty=10000\arraybackslash}X}
\toprule
Model & Steps & Whole answer, chain (\%) \mbox{[95\% interval]}, original rule & Whole answer, pieces (\%) \mbox{[95\% interval]}, original rule & Whole answer, pooled (\%) \mbox{[95\% interval]}, original rule & Format compliance, chain (\%), original rule & Format flag fires, original rule \\
\midrule
Haiku-medium & 4 & 42.00~\mbox{[26.50, 59.50]} & 69.00~\mbox{[54.00, 82.00]} & 55.50~\mbox{[41.50, 68.25]} & 99.00 & no \\
\addlinespace
Haiku-medium & 6 & 46.00~\mbox{[30.50, 61.50]} & 50.00~\mbox{[34.50, 65.50]} & 48.00~\mbox{[33.75, 62.00]} & 98.00 & no \\
\addlinespace
Haiku-medium & 8 & 54.00~\mbox{[39.50, 68.00]} & 50.50~\mbox{[39.00, 61.50]} & 52.25~\mbox{[42.00, 63.25]} & 98.50 & no \\
\addlinespace
Haiku-medium & 12 & 23.50~\mbox{[13.00, 36.00]} & 30.50~\mbox{[15.00, 47.00]} & 27.00~\mbox{[17.00, 38.25]} & 96.00 & no \\
\addlinespace
Haiku-medium & 16 & 14.50~\mbox{[6.50, 24.50]} & 21.50~\mbox{[13.00, 30.50]} & 18.00~\mbox{[10.25, 26.00]} & 94.00 & yes \\
\addlinespace[1em]
Opus-medium & 4 & 95.00~\mbox{[76.39, 99.11]} & 100.00~\mbox{[83.89, 100.00]} & 97.50~\mbox{[92.50, 100.00]} & 95.00 & yes \\
\addlinespace
Opus-medium & 6 & 100.00~\mbox{[83.89, 100.00]} & 100.00~\mbox{[83.89, 100.00]} & 100.00~\mbox{[83.16, 100.00]} (exact, post~hoc; 20~puzzles) & 100.00 & no \\
\addlinespace
Opus-medium & 8 & 95.00~\mbox{[76.39, 99.11]} & 100.00~\mbox{[83.89, 100.00]} & 97.50~\mbox{[92.50, 100.00]} & 100.00 & no \\
\addlinespace
Opus-medium & 12 & 95.00~\mbox{[76.39, 99.11]} & 100.00~\mbox{[83.89, 100.00]} & 97.50~\mbox{[92.50, 100.00]} & 95.00 & yes \\
\addlinespace
Opus-medium & 16 & 100.00~\mbox{[83.89, 100.00]} & 100.00~\mbox{[83.89, 100.00]} & 100.00~\mbox{[83.16, 100.00]} (exact, post~hoc; 20~puzzles) & 100.00 & no \\
\addlinespace[1em]
Sonnet-medium & 4 & 55.00~\mbox{[34.21, 74.18]} & 95.00~\mbox{[76.39, 99.11]} & 75.00~\mbox{[62.50, 87.50]} & 60.00 & yes \\
\addlinespace
Sonnet-medium & 6 & 70.00~\mbox{[48.10, 85.45]} & 75.00~\mbox{[53.13, 88.81]} & 72.50~\mbox{[57.50, 87.50]} & 75.00 & yes \\
\addlinespace
Sonnet-medium & 8 & 75.00~\mbox{[53.13, 88.81]} & 65.00~\mbox{[43.29, 81.88]} & 70.00~\mbox{[57.50, 82.50]} & 90.00 & yes \\
\addlinespace
Sonnet-medium & 12 & 95.00~\mbox{[76.39, 99.11]} & 50.00~\mbox{[29.93, 70.07]} & 72.50~\mbox{[60.00, 85.00]} & 100.00 & no \\
\addlinespace
Sonnet-medium & 16 & 75.00~\mbox{[53.13, 88.81]} & 55.00~\mbox{[34.21, 74.18]} & 65.00~\mbox{[50.00, 77.50]} & 100.00 & no \\
\bottomrule
\end{tabularx}
% note: Whole-answer accuracy is the share of samples with every step right; a chain is right only if every step line and the answer line are right. Both marking rules are shown side by side: the Deviation~2 rule, which is official, and the original rule, under which a step line written more than once fails even when the repeats are identical.
% note: Format compliance sits beside the whole-answer accuracy: the share of answered chain samples whose format was ok under each marking rule (the files give it for chains only), with the format flag, which fires when more than 5\% of the answered chain samples are not ok. Panel~3 adds the chain samples answered and not placeable.
% note: Three flags sit at exactly 5.0\% non-compliance; the locked check raises them through floating-point rounding, though its rule reads `more than 5\%'; we report them as stamped (Opus-medium's two cells sit at exactly 5.0\% under the original rule's columns too, raised the same way).
% note: Per-step accuracy (Panel~1): for pieces, the share of a sample's pieces that are right (a missing or unanswered sample scores nothing); for chains, each step is judged from the model's own previous step, and only samples whose step lines could be placed count: unplaceable chain samples are left out (Panel~3 counts them).
% note: Each length has its own small set: 20 puzzles. Haiku-medium has 10 samples per puzzle; Opus-medium and Sonnet-medium have 1.
% note: Intervals are 95\% and unadjusted: cluster bootstrap over puzzles with 2,000 resamples, except Opus-medium and Sonnet-medium per-shape whole-answer intervals, which are Wilson (the file's interval field).
% note: ``(exact, post hoc; \ldots{} puzzles)'' marks a bootstrap interval that collapsed to one value (every puzzle right, or every one wrong), replaced by the exact two-sided 95\% Clopper--Pearson interval from free analysis 4 (post hoc; it changes no test). The count in the mark is the number of puzzles, not of samples (the conservative unit).
% note: In Opus-medium's per-step chain cells the count in the mark is only the puzzles whose chain sample could be placed: the set's puzzles less the unplaceable chain samples (Panel~3).
% note: ``(zero width)'' marks Opus-medium's $q$ interval (Panel~4) and $q^{N}$ bands (Panel~5), which collapsed to one value; free analysis 4 gives them no exact interval (post hoc): ``$q$ pools every piece-step over all lengths and sets (clusters stratified by set), so no single count of base puzzles carries it, and its $q^{N}$ bands are powers of its interval''. Its note on $q$ adds: ``the length sweep's $q$ is estimated from pieces only, so its exact counterpart is the pieces per-step row''; Opus-medium's pieces per-step cells carry exact intervals at every length under both marking rules (Panel~1). In the first quotation the file's ``base puzzles'' are the paper's puzzles.
% note: Haiku-medium's vote (Panel~2) is over all 10 of its samples per puzzle and is not cost-matched: its cost is not set equal to that of another model's answer.
% note: $q$ is each model's share of pieces right, estimated from the pieces puzzles only and pooled over the lengths; $q^{N}$ is the independent-pieces reference and observed / $q^{N}$ compares the observed whole-answer accuracy with it. The 24-step row is an extrapolation only; its label, quoted as stamped, says ``dotted'', and the figures draw the line dashed.
% note: Pilot~1 markers are at 24 steps, exploratory, and ran before Lock~2 with a different rule mix, different sets and different timing; they never enter $q$. A preset is one setting of the puzzle generator (its rule kinds and dials). One marker pools Pilot~1's sets of three rule-set presets; the other only those of the one preset used later, whose rules the length sweep keeps at every length.
% note: The sweep is descriptive (pre-registered as such); no test is run on it.
